\documentclass[11pt]{article}
\usepackage[T1]{fontenc}
\usepackage{lmodern}
\usepackage[margin=1in]{geometry}
\usepackage{amsmath,amssymb,amsthm,booktabs,microtype}
\usepackage[hidelinks]{hyperref}
\newtheorem{theorem}{Theorem}
\newtheorem{proposition}[theorem]{Proposition}
\newcommand{\Z}{\mathbb Z}
\newcommand{\Q}{\mathbb Q}
\newcommand{\ord}{\operatorname{ord}_z}
\title{An explicit formal quotient for\\Class 214 inversion sequences}
\author{Xiaojian Zeng\\[3pt]
  \small PBC School of Finance, Tsinghua University, 100190\\
  \small\texttt{xiaojianzeng@foxmail.com}}
\date{}
\begin{document}
\maketitle

\begin{abstract}
We determine the ordinary generating function for inversion sequences
with no indices $i<j<k$ satisfying both $e_i\ge e_k$ and $e_j\ge e_k$,
the avoidance class known as Class 214. Starting from the auxiliary
functional equation of Britt and Beaton, we give an explicit quotient
of two coefficientwise convergent sums, with no unknown counting or
boundary series. Subtracting the constant and linear terms of a boundary
series makes its iterated remainder vanish in the formal topology.
Truncation to $m$ summands gives an algebraic series $F_m$ with the exact
contact $F_m-F=z^{m+3}+O(z^{m+4})$ for every $m\ge0$.
This is a generating-function representation result: the original
non-algebraicity conjecture remains unresolved.
\end{abstract}

\section{The counting problem and the quotient}

Let $\mathcal I_n$ consist of the integer sequences $(e_1,\ldots,e_n)$
such that $0\le e_i<i$ and
\begin{equation}\label{eq:avoid}
  \text{there are no }i<j<k\text{ with }e_i\ge e_k\text{ and }e_j\ge e_k.
\end{equation}
The empty sequence is counted once. Write
$a_n=|\mathcal I_n|$ and $F(z)=\sum_{n\ge0}a_nz^n$.
Equivalently, the forbidden subsequence patterns are
$000,010,100,110,120,210$, where a pattern records all strict-order
and equality relations, without an adjacency requirement.
In the relation-triple classification of Martinez and Savage
\cite{MartinezSavage}, this is the class $(-,\ge,\ge)$, labeled 214.
We use the relation condition \eqref{eq:avoid}; the correction to its
older pattern-list annotation is recorded in Britt and Beaton's Table~1
\cite{BrittBeaton}.

Britt and Beaton \cite[Section~3.5, equations~(3.31)--(3.33)]{BrittBeaton}
give a succession rule and an auxiliary functional equation for this class,
and leave the generating function unsolved there. Mansour
\cite[Table~2, class~44, and Theorem~13]{Mansour} places the same six-pattern
set in a Wilf class of five sets and gives generating-tree descriptions.
Those descriptions and equivalences are prior results. The contribution
here is an explicit evaluation of the remaining boundary coefficient,
together with a sharp error formula for finite truncations.
The proof uses kernel substitution, a standard method illustrated by
Prodinger \cite{Prodinger}, followed by a regularized formal iteration.

The original question also asks whether $F$ is non-algebraic over $\Q(z)$:
whether $P(z,F(z))\ne0$ for every nonzero $P\in\Q[z,w]$.
We do not decide that question. All limits below are formal,
coefficientwise limits at $z=0$; $O(z^r)$ denotes an element of
$z^r\Z[[z]]$ when the coefficients under discussion are integral.

Define the small root
\begin{equation}\label{eq:X}
 X=\frac{1-2z-\sqrt{1-4z}}{2z}
   =z+2z^2+5z^3+\cdots,\qquad X=z(1+X)^2,
\end{equation}
where $\sqrt{1-4z}=1+O(z)$. Set
\begin{equation}\label{eq:qb}
 q=\frac{X^2}{1+2X},\qquad
 b=\frac{X^3}{(1+X)(1+2X)},
\end{equation}
and introduce the rational functions
\begin{align}
 T(u)&=\frac{zu}{1-zu}, & H(u)&=z(1+u)^2-u,\label{eq:TH}\\
 R_0(u)&=-\frac{zu^2(1-z-zu)}{1-zu}, &
 R_1(u)&=\frac{u^2(1-z-zu)^2}{1-zu}.\label{eq:R}
\end{align}
For $j\ge0$, put
\begin{equation}\label{eq:orbit}
 t_j=\frac{z^jq}{1-q\sum_{r=1}^jz^r},\qquad
 D_j=\prod_{i=0}^jH(t_i),\qquad D_{-1}=1.
\end{equation}
Empty sums are zero throughout.

\begin{theorem}\label{thm:quotient}
The sums
\begin{equation}\label{eq:sums}
 A=\sum_{j\ge0}\frac{R_0(t_j)}{D_j},\qquad
 L=q+\sum_{j\ge0}\frac{R_1(t_j)}{D_j}
\end{equation}
converge coefficientwise in $\Z[[z]]$. Their displayed summands have
exact orders $j+4$ and $j+3$, respectively. Moreover,
\begin{equation}\label{eq:quotient}
 F(z)=\frac{b-A}{zL},\qquad
 L=z^2+O(z^3),\qquad b-A=z^3+O(z^4).
\end{equation}
For every integer $m\ge0$, define
\begin{equation}\label{eq:finite}
 A_m=\sum_{j=0}^{m-1}\frac{R_0(t_j)}{D_j},\qquad
 L_m=q+\sum_{j=0}^{m-1}\frac{R_1(t_j)}{D_j},\qquad
 F_m=\frac{b-A_m}{zL_m}.
\end{equation}
Then $F_m\in\Q(z,X)\cap\Z[[z]]$, and
\begin{equation}\label{eq:contact}
 F_m-F=z^{m+3}+O(z^{m+4}).
\end{equation}
\end{theorem}

Thus $m=\max(0,N-2)$ summands suffice to determine the coefficients
through degree $N$. This assertion concerns exact coefficient contact,
not a running-time comparison with the published recurrence.
For example, $F_0=b/(zq)=1+X$, whose coefficient of $z^3$ exceeds that of
$F$ by one. The initial series is
\[
 F=1+z+2z^2+4z^3+10z^4+26z^5+73z^6+214z^7+651z^8+O(z^9).
\]

\section{The inherited auxiliary equation}

We recall the combinatorial input of \cite[Section~3.5]{BrittBeaton}.
Enlarge the class to $\mathcal B$ by allowing occurrences of $000,010,100$
only when every entry represented by a $0$ is literally zero.
The patterns $110,120,210$ remain forbidden. There are at most two
runs of zeros: three runs would provide two positive entries before
a later zero, creating one of these last three patterns.
Let $p,s$ be the lengths of the first and second zero runs, taking
$s=0$ when the second is absent and $(p,s)=(0,0)$ for the empty object.

For completeness, the published left-growth rule has root $(0,0)$ and
the following children, with one child for each listed choice:
\begin{equation}\label{eq:rule}
\begin{aligned}
 (p,s)&\longmapsto(p+1,s),\\
 (p,s)&\longmapsto(p+1,s-1) &&(s>0),\\
 (p,0)&\longmapsto(p-\ell,\ell) &&(0\le\ell\le p-1),\\
 (p,0)&\longmapsto(p-\ell-1,\ell) &&(0\le\ell\le p-2).
\end{aligned}
\end{equation}
This operation increments the positive entries, promotes selected
zeros to ones, and prepends a zero. If $s=0$, one may promote one zero,
or two provided one is rightmost. If $s>0$, only the rightmost zero
can be promoted. Deleting the initial zero and decrementing positive
entries recovers the parent. These are the choices encoded in
\eqref{eq:rule}.

Let $B(z,x,y)$ mark length, $p$, and $s$, and abbreviate
$B(x,y)=B(z,x,y)$ and $V(x)=B(z,x,0)$.
Every coefficient in $z$ is an integer polynomial in the marking
variables. The desired class is exactly $s=0$, $p\le2$.
Indeed a second zero run produces $010$, and three initial zeros
produce $000$; a single run of at most two zeros leaves none of the
allowed exceptions. Consequently, if $c=[x]V$ and $d=[x^2]V$, then
\begin{equation}\label{eq:extract}
 F=1+c+d,\qquad
 V(x)=1+zx+z^2(x+x^2)+O(z^3).
\end{equation}

Multiplying equation~(3.33) of \cite{BrittBeaton} by $y-x$ gives
\begin{equation}\label{eq:cleared}
\begin{split}
0={}&(y+z)(y-x)-(y-zx-zxy)(y-x)B(x,y)\\
 &+xz(1+y)V(y)-z(y+2xy-x^2)V(x).
\end{split}
\end{equation}
Set $y=T(x)$. This substitution is valid in $\Q[x][[z]]$,
since $1-zx$ is a unit and $T(x)$ is divisible by $z$.
The coefficient of $B$ vanishes. Rational simplification gives
\[
 \frac{T+2xT-x^2}{x(1+T)}=H(x),\qquad
 -\frac{(T+z)(T-x)}{xz(1+T)}
 =S(x):=\frac{(1+x-zx)(1-z-zx)}{1-zx}.
\]
Canceling the nonzero common factor $xz(1+T)$ in this integral domain
yields the boundary equation
\begin{equation}\label{eq:boundary}
 H(x)V(x)=V(T(x))-S(x).
\end{equation}
We have cleared the denominator before substitution; no evaluation at
a pole is involved.

Since
\[
 S(x)=1-z+(1-2z)x-\sum_{k\ge2}(zx)^k,
\]
the coefficient of $x^2$ in \eqref{eq:boundary} is
\[
 zd+(2z-1)c+z=z^2(c+d+1).
\]
Cancellation of the unit $1-z$ gives $zd=(1-z)c-z$, so
\begin{equation}\label{eq:c}
 F=c/z.
\end{equation}
This is an ordinary formal series because $c=z+O(z^2)$.

\section{Regularization and the exact remainder}

Substitute $x=X$ in the undivided equation \eqref{eq:boundary}.
As $H(X)=0$, we obtain $V(T(X))=S(X)$.
Replacing $z$ by $X/(1+X)^2$ in the rational expressions gives
\begin{equation}\label{eq:kernelvalue}
 T(X)=q,\qquad S(X)=1+b,\qquad V(q)=1+b.
\end{equation}
In particular $q=z^2+O(z^3)$ and $b=z^3+O(z^4)$.
The coefficients of $X$ are integral by the recursion in \eqref{eq:X}.

Set $W(x)=V(x)-1-cx$. Equation~\eqref{eq:extract} implies the precise
initial form
\begin{equation}\label{eq:Winitial}
 W(x)=z^2x^2+z^3x^2U(z,x),\qquad U\in\Z[x][[z]].
\end{equation}
Subtracting the constant and linear parts of \eqref{eq:boundary} gives
\begin{equation}\label{eq:regular}
 H(x)W(x)=W(T(x))+R_0(x)+cR_1(x),
\end{equation}
where direct calculation gives
$R_0=1-S-H$ and $R_1=T-Hx$, exactly as in \eqref{eq:R}.

The denominator in \eqref{eq:orbit} shows directly that
$t_0=q$ and $t_{j+1}=T(t_j)$. For every $j\ge0$,
\begin{align}
 t_j&=z^{j+2}(1+O(z)), &
 H(t_j)&=z(1+O(z)),\nonumber\\
 D_j&=z^{j+1}(1+O(z)), &
 R_0(t_j)&=-z^{2j+5}(1+O(z)),\label{eq:orders}\\
 &&R_1(t_j)&=z^{2j+4}(1+O(z)).\nonumber
\end{align}
All units occurring here have integer coefficients and constant term
one. Their inverses are integral as well. Thus the quotients
$R_0(t_j)/D_j$ and $R_1(t_j)/D_j$ belong to $\Z[[z]]$, have exact
orders $j+4$ and $j+3$, and have leading coefficients $-1$ and $1$.
These orders tend to infinity, proving convergence of both sums
in \eqref{eq:sums}.

Iterating \eqref{eq:regular} finitely, including the empty iteration
when $m=0$, gives
\begin{equation}\label{eq:iteration}
 W(q)=\sum_{j=0}^{m-1}\frac{R_0(t_j)+cR_1(t_j)}{D_j}
       +\frac{W(t_m)}{D_{m-1}}.
\end{equation}
By \eqref{eq:Winitial} and \eqref{eq:orders}, its remainder satisfies
\begin{equation}\label{eq:tail}
 W(t_m)=z^{2m+6}(1+O(z)),\qquad
 \frac{W(t_m)}{D_{m-1}}=z^{m+6}(1+O(z)).
\end{equation}
The term $z^3x^2U$ contributes only at larger order after substitution,
so the leading coefficient in \eqref{eq:tail} is exactly one.
This proves that the remainder vanishes in the formal limit.
Without subtracting the constant and linear terms of $V$, the iterated
remainder would not vanish: $V(t_m)/D_{m-1}$ has order $-m$.

Equation~\eqref{eq:kernelvalue} now gives
\[
 b=cq+W(q)=A+cL.
\]
Together with \eqref{eq:c}, this proves \eqref{eq:quotient}.
The orders already established imply $L=z^2+O(z^3)$ and
$b-A=z^3+O(z^4)$. Dividing both numerator and denominator by $z^3$
therefore makes \eqref{eq:quotient} a quotient of integral units.

Keeping the finite remainder instead gives
\[
 b=A_m+cL_m+\frac{W(t_m)}{D_{m-1}},\qquad
 F_m-F=\frac{W(t_m)}{D_{m-1}zL_m}.
\]
For every $m\ge0$, $L_m=z^2+O(z^3)$; hence \eqref{eq:tail} proves
\eqref{eq:contact}, including its sign and the case $m=0$.
Similarly $b-A_m=z^3+O(z^4)$, so $F_m\in\Z[[z]]$.
Every finite sum and product in its definition lies in $\Q(z,X)$.
This proves all assertions of Theorem~\ref{thm:quotient}.

\begin{proposition}\label{prop:unique}
Equation~\eqref{eq:boundary} has at most one solution in $\Q[x][[z]]$.
\end{proposition}
\begin{proof}
For the difference $U$ of two solutions, $H(x)U(x)=U(T(x))$.
Setting $x=0$ gives $(z-1)U(0)=0$, hence every coefficient polynomial
of $U$ is divisible by $x$. If $z^ku_k(x)$ is the first nonzero term
of $U$, then $U(T(x))$ has order at least $k+1$, whereas $H(x)U(x)$
has nonzero coefficient $-xu_k(x)$ at degree $k$. This is impossible.
\end{proof}
The combinatorial series supplies existence. No additional degree-growth
assumption on the coefficient polynomials is used in this uniqueness
argument.

\section{Exact checks and limitations}

The proof above is an all-orders formal argument. Exact computations
provide separate implementation checks. Direct expansion of the quotient,
the published left-growth rule, and a right-append dynamic program agree
through degree $40$. A separately implemented check evaluates the quotient
in the parameter $w=X$, with $z=w/(1+w)^2$, and agrees with counting
through degree $48$. It uses rational polynomial arithmetic with exclusive
working precision $120$ in $w$. Lagrange inversion
\cite[Theorem~2.1.1]{Gessel} converts back by
\begin{equation}\label{eq:lagrange}
 [z^n]f(X(z))=\frac1n\sum_{k=1}^n
 k[w^k]f(w)\binom{2n}{n-k}\quad(n\ge1).
\end{equation}
In this verification $f(w)$ is the expression after substituting
$z=w/(1+w)^2$, so it has no remaining independent $z$ argument.

The counting criterion is particularly short. An entry $v$ appended to
a valid prefix is allowed exactly when at most one earlier entry is
at least $v$. A dynamic program therefore retains the two largest entries
with multiplicity, padded by $-1$ initially, and requires $v$ to exceed
the smaller retained entry. The inversion-sequence bound at the new
position must also be imposed. Literal tests of indexed triples on all
inversion sequences through length $9$ give the same counts.
The parameter-based verifier also checks the first error in
\eqref{eq:contact} for $0\le m\le45$ and the summand orders for
$0\le j\le47$. These are finite checks, not proofs of the quantified
statements. The accompanying reproducibility instructions identify the
programs and their exact coverage.

The formula removes the unknown boundary series from the inherited
equation. No equivalent solved quotient was located in the inspected
primary sources and the bounded literature searches through
13 September 2026; this is not an exhaustive priority assertion.
The infinite limit need not belong to the algebraic field containing
every $F_m$. Neither the algebraicity of the approximants nor individual
zeros of their denominators decide the algebraic status or analytic
singularities of $F$: cancellations in the infinite quotient remain to
be controlled. We make no analytic-continuation or asymptotic claim.
The non-algebraicity part of the original problem remains unresolved.

An obstruction to a prescribed coefficient formula has a different scope.
Z.~Zeng gives such a two-basis obstruction for phylogenetic-network counts
\cite{ZengSquareRoot}.
The local author manuscript uses the nonzero coefficient $[Y^3]R_4(u)=5/16$, with
$Y=\sqrt{1-2u}$ and $R_4$ the exponential generating function for
reticulation-visible networks with four reticulations, to exclude a
representation of the difference between
reticulation-visible and tree-child counts by polynomial multiples of
$(2\ell-3)!!$ and $2^{\ell-3}\ell!$ for all $\ell\ge5$.
The obstruction concerns these two specified bases; it supplies no
non-algebraicity criterion for the present quotient.
The cited version is the dated local author manuscript identified in
the bibliography; see Theorem~1.1, Corollary~3.2, and Lemma~4.1 there.

\paragraph{Computational and AI assistance.}
OpenAI language models assisted with the derivation, exact Python
implementations, literature searches, mathematical checking, and manuscript
preparation. A separate model context reconstructed the argument and
implemented the check in the parameter $w$ using SymPy~1.14.0.
That additional review is not human peer review or a claim of statistically
independent evidence. The computations use integer or rational arithmetic;
no floating-point fit or formal proof-assistant certification is used.

\bibliographystyle{plain}
\bibliography{references}
\end{document}
